% from ega4-2-fr.tex
\oldpage[IV-2]{11}
On notera que cette notion est locale sur $Y$, mais \emph{non sur} $X$~; en
particulier on peut avoir $\mathscr F_x=0$ pour certains $x\in X$, autrement
dit $\operatorname{Supp}(\mathscr F)$ n'est pas nécessairement égal à $X$.
Toutefois, il résulte aussitôt du critère (2.2.1, \emph{b)}) que pour tout
$y\in Y$, il existe au moins un $x\in f^{-1}(y)$ tel que
$(\mathscr F_y)_x=\mathscr F_x\otimes_{\mathcal O_y}k(y)\neq0$, et
\emph{a fortiori} $\mathscr F_x\neq0$~; en d'autres termes~:

\begin{corollary}[2.2.5]
\label{IV.2.2.5-fr}
Si $\mathscr F$ est un $\mathcal O_X$-Module quasi-cohérent, fidèlement plat
relativement à $f$, on a $f(\operatorname{Supp}(\mathscr F))=Y$, et
\emph{a fortiori} $f$ est un morphisme surjectif.
\end{corollary}

Ce résultat admet une réciproque partielle~:

\begin{corollary}[2.2.6]
\label{IV.2.2.6-fr}
Soit $\mathscr F$ un $\mathcal O_X$-Module quasi-cohérent de type fini. Pour
que $\mathscr F$ soit fidèlement plat relativement à $f$, il faut et il
suffit que $\mathscr F$ soit $f$-plat et que
$f(\operatorname{Supp}(\mathscr F))=Y$.
\end{corollary}

En effet (I, 9.1.13 et 3.6.1) on a
$\operatorname{Supp}(\mathscr F_y)=f^{-1}(y)\cap\operatorname{Supp}(\mathscr F)$,
donc dans ce cas le critère (2.2.1, \emph{b)}) n'est autre que la condition du
corollaire. En particulier, le $\mathcal O_X$-Module $\mathcal O_X$ est
fidèlement plat relativement à $f$ si et seulement s'il est $f$-plat et si
$f$ est surjectif, autrement dit si et seulement si \emph{le morphisme $f$
est fidèlement plat} $(0_{\mathrm I},6.7.8)$.

Explicitons les propriétés intervenant dans la définition (2.2.4)~:

\begin{proposition}[2.2.7]
\label{IV.2.2.7-fr}
Soient $f:X\to Y$ un morphisme de préschémas, $\mathscr F$ un
$\mathcal O_X$-Module quasi-cohérent fidèlement plat relativement à $f$.
Alors, pour qu'une suite $\mathscr G'\to\mathscr G\to\mathscr G''$ de
$\mathcal O_Y$-Modules quasi-cohérents soit exacte, il faut et il suffit que
la suite correspondante
$\mathscr F\otimes_Y\mathscr G'\to\mathscr F\otimes_Y\mathscr G
\to\mathscr F\otimes_Y\mathscr G''$ le soit. En particulier, pour qu'un
homomorphisme $u:\mathscr G\to\mathscr G'$ de $\mathcal O_Y$-Modules
quasi-cohérents soit injectif (resp. surjectif, bijectif, nul), il faut et il
suffit que $1_{\mathscr F}\otimes f^*(u):
\mathscr F\otimes_Y\mathscr G\to\mathscr F\otimes_Y\mathscr G'$ le soit.
Pour qu'un $\mathcal O_Y$-Module quasi-cohérent $\mathscr G$ soit nul, il faut
et il suffit que $\mathscr F\otimes_Y\mathscr G=0$. Pour tout
$\mathcal O_Y$-Module quasi-cohérent $\mathscr G$, l'application
$\mathscr G'\mapsto\mathscr F\otimes_Y\mathscr G'$ de l'ensemble des
sous-$\mathcal O_Y$-Modules quasi-cohérents de $\mathscr G$ dans l'ensemble
des sous-$\mathcal O_X$-Modules quasi-cohérents de
$\mathscr F\otimes_Y\mathscr G$ est injective.
\end{proposition}

Pour démontrer la dernière assertion, c'est-à-dire que pour deux
sous-$\mathcal O_Y$-Modules quasi-cohérents $\mathscr G'$, $\mathscr G''$ de
$\mathscr G$, la relation
$\mathscr F\otimes_Y\mathscr G'=\mathscr F\otimes_Y\mathscr G''$ entraîne
$\mathscr G'=\mathscr G''$, on peut (en remplaçant $\mathscr G''$ par
$\mathscr G'+\mathscr G''$) se ramener au cas où
$\mathscr G'\subset\mathscr G''$, et il suffit alors d'appliquer la seconde
assertion de l'énoncé à l'injection $u:\mathscr G'\to\mathscr G''$.

\begin{corollary}[2.2.8]
\label{IV.2.2.8-fr}
Soit $f:X\to Y$ un morphisme fidèlement plat. Pour tout
$\mathcal O_Y$-Module quasi-cohérent $\mathscr G$, l'application canonique
\begin{equation}
\tag{2.2.8.1}
\label{IV.2.2.8.1-fr}
\Gamma(Y,\mathscr G)\longrightarrow\Gamma(X,f^*(\mathscr G))
\end{equation}
est injective.
\end{corollary}

En effet, $\Gamma(Y,\mathscr G)$ s'identifie canoniquement à
$\operatorname{Hom}_{\mathcal O_Y}(\mathcal O_Y,\mathscr G)$
$(0_{\mathrm I},5.1.1)$, et de même $\Gamma(X,f^*(\mathscr G))$ à
$\operatorname{Hom}_{\mathcal O_X}(\mathcal O_X,f^*(\mathscr G))$. En vertu
de (2.2.1) et (2.2.4) l'hypothèse entraîne que le foncteur
$\mathscr G\mapsto f^*(\mathscr G)$ est exact et fidèle dans la catégorie des
$\mathcal O_Y$-Modules quasi-cohérents, et par suite un homomorphisme
$u:\mathcal O_Y\to\mathscr G$ est nul si et seulement si
$f^*(u):f^*(\mathcal O_Y)=\mathcal O_X\to f^*(\mathscr G)$ est nul.

\begin{remark}[2.2.9]
\label{IV.2.2.9-fr}
Les résultats de (2.2.7) et (2.2.8) sont encore vrais lorsque les
$\mathcal O_Y$-Modules $\mathscr G'$, $\mathscr G$, $\mathscr G''$ qui y
figurent sont quelconques (non nécessairement quasi-cohérents). En effet,
pour tout $y\in Y$, il existe $x\in f^{-1}(y)$ tel que $\mathscr F_x$ soit un
$\mathcal O_y$-module
